\documentclass[10pt,a4paper]{amsart}
\usepackage[margin=19mm]{geometry}
\usepackage{amsmath,amssymb,mathtools}
\usepackage[scr=rsfs,cal=euler]{mathalfa}
\usepackage{xfrac}
\usepackage[unicode,hidelinks,bookmarks=false]{hyperref}
\newcommand{\ie}{i.e.\ }
\newcommand{\cf}{cf.\ }
\newcommand{\resp}{respectively\ }
\newcommand{\Subsection}{\subsection}
\providecommand{\nobreakdash}{\nobreak\hbox{-}}
\setlength{\emergencystretch}{2em}
\multlinegap0pt
\usepackage[shortlabels]{enumitem}
\usepackage{amssymb}
\usepackage{mathtools}
\newtheorem{lemma}{Lemma}
\usepackage{cleveref}
\crefname{lemma}{Lemma}{Lemmas}
\title{GIT quotient of minimal dimensional Schubert variety modulo a subtorus}
\author{Arkadev Ghosh, S. S. Kannan}
\date{Source: arXiv:2604.25645v1 (CC BY 4.0)}
\begin{document}
\maketitle

\noindent\emph{Source and licence.} GIT quotient of minimal dimensional Schubert variety modulo a subtorus. Version arXiv:2604.25645v1 (CC BY 4.0), distributed by arXiv under the Creative Commons Attribution 4.0 International licence, http://creativecommons.org/licenses/by/4.0/, as recorded in the arXiv licence metadata for this exact version and shown on the abstract page https://arxiv.org/abs/2604.25645v1. Full licence notice: \texttt{licenses/arxiv-2604.25645-CC-BY-4.0.txt}.\par\smallskip

\noindent\textit{Editorial scope.} Counted unit: the article's lemma 9 (source0.tex lines 1045--1047). The article's complete original proof of it is delivered with it (source0.tex lines 1048--1065, one \texttt{proof} environment of 18 lines). No mathematics is written, repaired or re-derived: the statement and the proof are the article's own lines, in the article's own notation, and no retained line is substituted or re-flowed.  Retained uncounted as the source-local context the proof needs: lines 459--461; lines 874--895; lines 956--962; lines 1019--1023; lines 1031--1033.  Nothing was omitted from any retained span, so the package carries no omission record.  No retained line contains a citation, so the wrapper carries no bibliography and no bibliography entry.  No retained line contains a figure environment, an image-inclusion command or drawing code, so no image is delivered and the package declares no asset dependency.  Editorial formatting only, in addition: an AMS \texttt{amsart} preamble that restates the article's own declarations and macros used by the retained lines, namely the environments \texttt{lemma} on separate counters, together with the standard packages those lines need.  The retained environments print in this document as Lemma 1 in order of appearance, whereas the article prints the same items with the numbering of its own full document.  The complete native source of the article as distributed in the arXiv e-print for this exact version is bundled unchanged as \texttt{sources/199317/source0.tex}.
For every $1\leq j\leq r$, define \begin{align}\label{eq 1.14}
	\gamma_{j}:=\sum_{k=0}^{m(i_{j},j)-1}\beta_{i_{e^{k}(i_{j},j)},e^{k}(i_{j},j)}
\end{align} 

\begin{lemma}\label{section:quotient:lemma 7}Let $J=(i_{1},i_{2},\cdots,i_{r})\in\Pi_{j=1}^{r}C_{j}$. 
	\begin{enumerate}
		\item 
		For $1\leq j\leq r-1$, 
		let $\delta_{J(j)}:=\sum_{k=0}^{m(i_{j},j)-1}\gamma_{i_{e^{k}(i_{j},j)},e^{k}(i_{j},j)}$. Then we have 
		\begin{align*}
			\langle\delta_{J(j)},\lambda^{\prime}_{lq}\rangle=\begin{dcases}
				1& \text{ if $1\leq l\leq j$}\\
				0& \text{ if $j+1\leq l\leq r-1$}
			\end{dcases}
		\end{align*}
		\item 	For $1\leq j\leq r$, 
		let $\gamma_{J(j)}:=\sum_{k=0}^{m(i_{j},j)-1}\beta_{i_{e^{k}(i_{j},j)},e^{k}(i_{j},j)}$. Then we have 
		\begin{align*}
			\langle\gamma_{J(j)},\lambda_{lq}\rangle=\begin{dcases}
				1& \text{ if $1\leq l\leq j$}\\
				0& \text{ if $j+1\leq l\leq r$}
			\end{dcases}
		\end{align*}
\text{$($Note that $\gamma_{J(j)}=\gamma_{j}$ as defined in \cref{eq 1.14}$)$}
	\end{enumerate}
\end{lemma}

	\begin{lemma}\label{section: quotient lemma 5}
		Let	$J=(i_{1},i_{2},\cdots,i_{r-1})\in \Pi_{j=1}^{r-1}C_{j}$. Then, we have 
		\begin{enumerate}
			\item $\tilde{U}(J)$ is stable under the action of $T^{\prime}_{J_{r-1}}$.
			\item $\tilde{V}(J)$ is stable under the action of $T_{J_{r}}$.
		\end{enumerate}
	\end{lemma}

	\begin{align}
		\text{For $1\leq j\leq r-1$, define\ }&
		b(J(j))(y):=\prod_{k=0}^{m(i_{j},j)-1}X_{\gamma_{i_{e^{k}(i_{j},j)},e^{k}(i_{j},j)}}(\tilde{u})\label{eq 5.1.12}\\
		&b(J(0))(y):=1\label{eq 5.1.14}
	\end{align}

	\begin{lemma}\label{section:quotient:lemma 8}
		Let $J\in\Pi_{j=1}^{r-1}C_{j}$. Then for every $t^{\prime}\in T^{\prime}_{J_{r-1}}$ and $1\leq j\leq r-1$,  we have $$b(J(j))(t^{\prime}y)=\delta_{J(j)}(t^{\prime})b(J(j))(y),\text{\ for all $y\in \tilde{U}(J)$}.$$
	\end{lemma}

	\begin{lemma}\label{section:quotient:lemma 9}
		Let $J_{1},J_{2}\in\Pi_{j=1}^{r-1}C_{j}$. Then for all $0\leq j\leq r-1$, $\tilde{b}_{J_{1},J_{2}}(j)$ is constant on $T^{\prime}_{J_{r-1}}$-orbits.
	\end{lemma}

	\begin{proof}
		Let $t^{\prime}\in T^{\prime}_{J_{r-1}}$. Let $t^{\prime}=\prod_{l=1}^{r-1}\lambda^{\prime}_{lq}(t_{l})$, where $t_{l}\in \mathbb{G}_{m}$ for all $1\leq l\leq r-1$. Let $y\in \tilde{U}(J_{1})\cap \tilde{U}(J_{2})$. From \cref{section: quotient lemma 5} $(1)$, $\tilde{U}(J_{1})\cap \tilde{U}(J_{2})$ is $T^{\prime}_{J_{r-1}}$-stable. Hence, $t^{\prime}y\in \tilde{U}(J_{1})\cap \tilde{U}(J_{2})$. Now for, $j=0$ we have $\tilde{b}_{J_{1},J_{2}}(0)(t^{\prime}y)=\frac{b(J_{1}(0))(t^{\prime}y)}{b(j_{2}(0))(t^{\prime}y)}=1$ (see \cref{eq 5.1.14}).\\
		Let $1\leq j\leq r-1$. Then, using \cref{section:quotient:lemma 8}, we have
		\begin{align}
			\tilde{b}_{J_{1},J_{2}}(j)(t^{\prime}y)&=\frac{b(J_{1}(j))(t^{\prime}y)}{b(J_{2}(j))(t^{\prime}y)}
			=\frac{\delta_{J_{1}(j)}(t^{\prime})b(J_{1}(j))(y)}{\delta_{J_{2}(j)}(t^{\prime})b(J_{2}(j))(y)}.\label{eq 5.1.39}
		\end{align}
	From \cref{section:quotient:lemma 7}(1), we have  \begin{align}
		\delta_{J_{1}(j)}(t^{\prime})&=\prod_{l=1}^{r-1}t_{l}^{\langle\delta_{J_{1}(j)},\lambda^{\prime}_{lq}\rangle}
		=\prod_{l=1}^{j}t_{l}\label{eq 5.1.40}
	\end{align}
Similarly, from \cref{section:quotient:lemma 7} (1), we have
\begin{align}
	\delta_{J_{2}(j)}(t^{\prime})=\prod_{l=1}^{j}t_{l}\label{eq 5.1.41}
\end{align}
		Therefore, \cref{eq 5.1.40}, \cref{eq 5.1.41}, we have $\delta_{J_{1}(j)}(t^{\prime})=\delta_{J_{2}(j)}(t^{\prime})$. Hence, 
		from \cref{eq 5.1.39}, we have 	$\tilde{b}_{J_{1},J_{2}}(j)(t^{\prime}y)=\tilde{b}_{J_{1},J_{2}}(j)(y)$.
	\end{proof}

\end{document}
